% Gleichungen lösen – bank/gleichungen-loesen/ (zusammenbau v0.4, Heft)
\einheitenkopf[t6]{Gleichungen lösen}
\setcounter{aufgabe}{27}

% Anlauf (ohne Prüfkennung)
\begin{aufgabe}{Schnittpunkte und Stellen – Anlauf}
\begin{teile}
\teil Welcher Weg? $x^4 - 10x^2 + 9 = 0$ \kreuz{Lösungsformel} \\ \kreuz{Substitution} \\ \kreuz{Polynomdivision}
\teil $f(x) = x^2 - 2x - 1$, $g(x) = x + 9$. Schnittpunkte? S1( \leerfeld | \leerfeld ), S2( \leerfeld | \leerfeld )
\end{teile}
\end{aufgabe}

\begin{aufgabe}{Schnittpunkte und Stellen – Prüfungsaufgaben}
\begin{teile}
\teil $f(x) = 0{,}5x^3 - 2x + 3$. An welchen Stellen schneidet der Graph die Gerade $y = 3$? \hfill (Abitur 2019 GK) \\ x1 = \leerfeld , x2 = \leerfeld , x3 = \leerfeld
\teil $f(x) = x^3 + 2x^2 - 8x - 1$. An welchen Stellen schneidet der Graph die Gerade $y = -1$? \hfill (Abitur 2019 GK) \\ x1 = \leerfeld , x2 = \leerfeld , x3 = \leerfeld
\teil Ein Mountainbiker springt von einer Rampe. Seine Flugbahn ist $f(x) = -0{,}02x^2 + 20$, das Gelände $g(x) = 0{,}0001x^4 - 0{,}05x^2 + 16$ ($x \ge 0$, 1 LE = 1 m). Koordinaten des Landepunkts $L$? \hfill (Abitur 2018 GK) \\ L( \leerfeld | \leerfeld )
\teil $f(x) = x^2 + 3x - 1$, $g(x) = -x^2 - x + 5$. Zeige, dass sich die Graphen nur für $x = -3$ und $x = 1$ schneiden. \hfill (Abitur 2019 GK)
\teil $f(x) = 0{,}5x^2 - x + 2$, $g(x) = -0{,}5x^2 + 3x + 2$. Zeige, dass sich die Graphen nur für $x = 0$ und $x = 4$ schneiden. \hfill (Abitur 2019 GK)
\teil $f(x) = 2x^2 - x - 7$, $g(x) = x^2 + 2x + 3$. Zeige, dass sich die Graphen nur für $x = -2$ und $x = 5$ schneiden. \hfill (Abitur 2019 GK)
\teil Ein Blatt in einer Grafik wird von $u(x) = 0{,}5x^3$ und $v(x) = x^2 \cdot (3 - x)$ begrenzt. Zeige, dass $P(0 | 0)$ und $Q(2 | 4)$ die einzigen gemeinsamen Punkte sind. \hfill (Abitur 2020 GK)
\teil Eine Flosse wird von $u(x) = 0{,}25x^3$ und $v(x) = 0{,}5x^2 \cdot (6 - x)$ begrenzt. Zeige, dass $P(0 | 0)$ und $Q(4 | 16)$ die einzigen gemeinsamen Punkte sind. \hfill (Abitur 2020 GK)
\teil Ein Segel wird von $u(x) = x^3$ und $v(x) = 2x^2 \cdot (3 - x)$ begrenzt. Zeige, dass $P(0 | 0)$ und $Q(2 | 8)$ die einzigen gemeinsamen Punkte sind. \hfill (Abitur 2020 GK)
\teil $h(x) = \frac{4}{x - 2}$ ($x \ne 2$), $g(x) = x - 2$. Koordinaten der gemeinsamen Punkte? \hfill (Abitur 2019 GK) \\ S1( \leerfeld | \leerfeld ), S2( \leerfeld | \leerfeld )
\teil $h(x) = \frac{9}{x + 4}$ ($x \ne -4$), $g(x) = x + 4$. Koordinaten der gemeinsamen Punkte? \hfill (Abitur 2019 GK) \\ S1( \leerfeld | \leerfeld ), S2( \leerfeld | \leerfeld )
\teil $h(x) = \frac{2}{x - 3}$ ($x \ne 3$), $g(x) = 0{,}5 \cdot (x - 3)$. Koordinaten der gemeinsamen Punkte? \hfill (Abitur 2019 GK) \\ S1( \leerfeld | \leerfeld ), S2( \leerfeld | \leerfeld )
\end{teile}
\end{aufgabe}

\begin{aufgabe}{Schnittpunkte und Stellen – Prüfungsaufgaben – weiter}
\begin{teile}
\teil Zeige: $-0{,}5 \cdot (x - 2) \cdot (x + 3)^2 = -0{,}5x^3 - 2x^2 + 1{,}5x + 9$. \hfill (Abitur 2020 GK)
\teil Zeige: $0{,}2 \cdot (x + 1) \cdot (x - 3)^2 = 0{,}2x^3 - x^2 + 0{,}6x + 1{,}8$. \hfill (Abitur 2020 GK)
\teil Zeige: $-2 \cdot (x + 4) \cdot (x - 1)^2 = -2x^3 - 4x^2 + 14x - 8$. \hfill (Abitur 2020 GK)
\end{teile}
\end{aufgabe}

% Anlauf (ohne Prüfkennung)
\begin{aufgabe}{Gleichungen mit e-Funktion und Logarithmus – Anlauf}
\begin{teile}
\teil Welcher Weg? $4e^{3x} - 20 = 0$ \kreuz{Logarithmieren} \\ \kreuz{e-Faktor kürzen} \\ \kreuz{e-Funktion anwenden}
\teil Löse: $4e^{2x} - 12 = 0$. Runde auf zwei Stellen. x = \leerfeld
\end{teile}
\end{aufgabe}

\begin{aufgabe}{Gleichungen mit e-Funktion und Logarithmus – Prüfungsaufgaben}
\begin{teile}
\teil $f(x) = (x^2 - 3x) \cdot e^{-0{,}5x}$, $g(x) = x \cdot e^{-0{,}5x}$ (1 LE = 1 dm). Wie weit liegen die Schnittstellen waagerecht auseinander? \hfill (Abitur 2021 GK) \\ \leerfeld[dm]
\teil $f(x) = (x - 4) \cdot e^{-x}$ hat die Ableitung $f'(x) = (5 - x) \cdot e^{-x}$. Zeige, dass sich die Graphen von $f$ und $f'$ bei $x = 4{,}5$ schneiden. \hfill (Abitur 2022 GK)
\teil $f(x) = (x^2 + 2x) \cdot e^{0{,}2x}$, $g(x) = 5x \cdot e^{0{,}2x}$ (1 LE = 1 dm). Wie weit liegen die Schnittstellen waagerecht auseinander? \hfill (Abitur 2021 GK) \\ \leerfeld[dm]
\teil $f(x) = e^{0{,}5x^2}$ hat die Ableitung $f'(x) = x \cdot f(x)$. Die Graphen von $f$ und $f'$ schneiden sich in genau einem Punkt. Steigung des Graphen von $f$ in diesem Punkt? \hfill (Abitur 2023 GK) \\ m = \leerfeld
\teil $f(x) = 2e^{0{,}25x^2}$ hat die Ableitung $f'(x) = 0{,}5x \cdot f(x)$. Die Graphen von $f$ und $f'$ schneiden sich in genau einem Punkt. Steigung des Graphen von $f$ in diesem Punkt? \hfill (Abitur 2023 GK) \\ m = \leerfeld
\end{teile}
\end{aufgabe}

\begin{aufgabe}{Schnittpunkte zweier Graphen durch Ausklammern eines gemeinsamen Faktors berechnen – Prüfungsaufgaben}
\begin{teile}
\teil $f(x) = (x - 2) \cdot e^{0{,}25x}$, $g(x) = x - 2$. Koordinaten der Schnittpunkte? \hfill (Abitur 2018 GK) \\ S( \leerfeld | \leerfeld ), T( \leerfeld | \leerfeld )
\teil $f(x) = (2x + 4) \cdot e^{-x}$, $g(x) = 2x + 4$. Koordinaten der Schnittpunkte? \hfill (Abitur 2018 GK) \\ S( \leerfeld | \leerfeld ), T( \leerfeld | \leerfeld )
\teil $f(x) = (3 - x) \cdot e^{x}$, $g(x) = 3 - x$. Koordinaten der Schnittpunkte? \hfill (Abitur 2018 GK) \\ S( \leerfeld | \leerfeld ), T( \leerfeld | \leerfeld )
\end{teile}
\end{aufgabe}

% Anlauf (ohne Prüfkennung)
\begin{aufgabe}{Gleichungen aufstellen und grafisch oder numerisch lösen – Anlauf}
\begin{teile}
\teil Welche Gleichung? Die Höhe $h(t)$ eines Ballons erreicht 50 m. \kreuz{$h(t) = 50$} \\ \kreuz{$h(50) = t$} \\ \kreuz{$h'(t) = 50$}
\teil Der Graph zeigt die Höhe $h(t)$ einer Drohne ($t$ in s, $h$ in m). Wann ist sie 4 m hoch? Schreibe die Gleichung und lies ab.

\begin{ksys}[xmin=0,xmax=6,ymin=0,ymax=6,ablesen] \parabel{-1}{3}{5}{h} \end{ksys}
Gleichung: \leerfeld , t = \leerfeld
\end{teile}
\end{aufgabe}

\begin{aufgabe}{Gleichungen aufstellen und grafisch oder numerisch lösen – Prüfungsaufgaben}
\begin{teile}
\teil Die Abbildung zeigt die Graphen von $h$ und $g$. Beschreibe, wie man die Lösungen von $h(x) - g(x) = 0$ am Bild findet, und lies alle ab. \hfill (Abitur 2019 GK) \\

\begin{ksys}[xmin=-4,xmax=4,ymin=-4,ymax=4,ablesen] \funktion{0.25*\x^3-2*\x}{h} \gerade{-1}{0}{g} \end{ksys}
x = \leerfeld
\teil Die Abbildung zeigt die Graphen von $p$ und $g$. Beschreibe, wie man die Lösungen von $p(x) - g(x) = 0$ am Bild findet, und lies alle ab. \hfill (Abitur 2019 GK) \\

\begin{ksys}[xmin=-2,xmax=5,ymin=-4,ymax=4,ablesen] \parabel{1}{1}{-3}{p} \gerade{1}{-2}{g} \end{ksys}
x = \leerfeld
\teil Die Abbildung zeigt die Graphen von $h$ und $g$. Beschreibe, wie man die Lösungen von $h(x) - g(x) = 0$ am Bild findet, und lies alle ab. \hfill (Abitur 2019 GK) \\

\begin{ksys}[xmin=-4,xmax=4,ymin=-4,ymax=4,ablesen] \funktion{0.25*\x^3-2*\x+1}{h} \gerade{-1}{1}{g} \end{ksys}
x = \leerfeld
\teil $a(x)$ ist die Zahl der Abonnenten eines Kanals in Hundert, $x$ in Wochen (Graph). Löse $a(x + 2) = a(x) + 6$ grafisch und deute die Lösung. \hfill (Abitur 2025 GK) \\

\begin{ksys}[xmin=0,xmax=6,ymin=0,ymax=10,ablesen] \parabel{0.5}{0}{0}{a} \end{ksys}
x = \leerfeld
\teil $a(x)$ ist die Zahl der Besucher einer Ausstellung in Tausend, $x$ in Tagen (Graph). Löse $a(x + 4) = a(x) + 6$ grafisch und deute die Lösung. \hfill (Abitur 2025 GK) \\

\begin{ksys}[xmin=0,xmax=6,ymin=0,ymax=8,ablesen] \parabel{0.25}{0}{1}{a} \end{ksys}
x = \leerfeld
\teil $w(t)$ ist die Wassermenge in einem Tank in m³, $t$ in Stunden (Graph). Löse $w(t + 2) = w(t) - 4$ grafisch und deute die Lösung. \hfill (Abitur 2025 GK) \\

\begin{ksys}[xmin=0,xmax=7,ymin=0,ymax=18,ystep=2,ablesen] \parabel{0.5}{6}{0}{w} \end{ksys}
t = \leerfeld
\end{teile}
\end{aufgabe}

\begin{aufgabe}{Gleichungen aufstellen und grafisch oder numerisch lösen – Prüfungsaufgaben – weiter}
\begin{teile}
\teil $f(x) = -x^2 + 6x - 4$. Bilde $f'(x)$ und löse $\frac{f(x) - 0}{x - 0} = f'(x)$ rechnerisch. \hfill (Abitur 2025 GK) \\ f'(x) = \leerfeld , x1 = \leerfeld , x2 = \leerfeld
\teil $f(x) = x^2 - 2x + 9$. Bilde $f'(x)$ und löse $\frac{f(x) - 0}{x - 0} = f'(x)$ rechnerisch. \hfill (Abitur 2025 GK) \\ f'(x) = \leerfeld , x1 = \leerfeld , x2 = \leerfeld
\teil $f(x) = -0{,}5x^2 + 3x - 8$. Bilde $f'(x)$ und löse $\frac{f(x) - 0}{x - 0} = f'(x)$ rechnerisch. \hfill (Abitur 2025 GK) \\ f'(x) = \leerfeld , x1 = \leerfeld , x2 = \leerfeld
\end{teile}
\end{aufgabe}

% Anlauf (ohne Prüfkennung)
\begin{aufgabe}{Ungleichungen – Anlauf}
\begin{teile}
\teil Welche Lösungen zählen? $x^2 \cdot (x - 3) > 0$ \kreuz{$x > 3$} \\ \kreuz{$x > 3$ und $x = 0$} \\ \kreuz{$x \ge 3$}
\teil Für welche $x$ gilt $(x - 1) \cdot (x + 2)^2 > 0$? \leerfeld
\end{teile}
\end{aufgabe}

\begin{aufgabe}{Ungleichungen – Prüfungsaufgaben}
\begin{teile}
\teil $f(x) = x^3 - 6x^2 + 11x - 1$, $g(x) = 0{,}5x^3 - 2x^2 + 3x - 1$. Bestimme rechnerisch alle $x$ mit $f(x) > g(x)$. \hfill (Abitur 2023 GK) \\ \leerfeld
\teil $f(x) = -x^3 + 5x^2 - 10x + 2$, $g(x) = -3x^2 + 6x + 2$. Bestimme rechnerisch alle $x$ mit $f(x) > g(x)$. \hfill (Abitur 2023 GK) \\ \leerfeld
\teil Die Düngerkonzentration im Boden ist $w_k(t) = k \cdot t \cdot e^{-0{,}1t} + 20$ ($t$ in Tagen, $k > 0$). Ab welchem $k$ beträgt sie am 15. Tag mindestens 40? \hfill (Abitur 2019 GK) \\ k = \leerfeld
\teil Der Pegel eines Stoffs in einem See ist $p_k(t) = k \cdot t \cdot e^{-0{,}2t} + 10$ ($t$ in Wochen, $k > 0$). Ab welchem $k$ beträgt er in Woche 10 mindestens 30? \hfill (Abitur 2019 GK) \\ k = \leerfeld
\teil Die Koffeinmenge im Blut ist $c_k(t) = k \cdot t \cdot e^{-0{,}5t} + 5$ ($t$ in h, $k > 0$). Ab welchem $k$ beträgt sie nach 4 Stunden mindestens 12? \hfill (Abitur 2019 GK) \\ k = \leerfeld
\teil $f(x) = 0{,}25x^4 - 2x^2 + 4$ wird für $x \ge 0$ durch $p(x) = -2x^2 + 4$ angenähert. Die Näherung ist brauchbar, solange der Betrag der Differenz höchstens $0{,}25$ ist. Beschreibe einen Lösungsweg, mit dem man alle $x \ge 0$ findet, für die $p$ brauchbar ist. \hfill (Abitur 2024 GK)
\teil $f(x) = x^4 - 4x^2 + 3$ wird für $x \ge 1$ durch $p(x) = -3x^2 + 3$ angenähert. Die Näherung ist brauchbar, solange der Betrag der Differenz höchstens $0{,}75$ ist. Beschreibe einen Lösungsweg, mit dem man alle $x \ge 1$ findet, für die $p$ brauchbar ist. \hfill (Abitur 2024 GK)
\end{teile}
\end{aufgabe}

\begin{aufgabe}{Intervall, in dem eine Modellfunktion mindestens einen vorgegebenen Wert annimmt, über eine Ungleichung bestimmen – Prüfungsaufgaben}
\begin{teile}
\teil Beim Erhitzen einer Flüssigkeit gilt $f(t) = 20 + 15t \cdot e^{-\frac{t}{8}}$ ($t$ in min, $f$ in °C). In welchem Zeitraum ist sie mindestens 60 °C heiß? Löse die Ungleichung mit dem Rechner (zwei Stellen). \hfill (Abitur 2017 GK) \\ von \leerfeld bis \leerfeld[min]
\teil Der senkrechte Abstand eines Gleitschirms zum Hang ist $d(x) = -0{,}0005x^4 + 0{,}1x^2 + 3$ ($0 \le x \le 14$, in m). In welchem Bereich beträgt er mindestens 6 m? \hfill (Abitur 2018 GK) \\ von \leerfeld bis \leerfeld[m]
\teil Die Konzentration eines Wirkstoffs ist $c(t) = 12t \cdot e^{-0{,}5t}$ ($t$ in h, $c$ in mg/l). Er wirkt ab 5 mg/l. In welchem Zeitraum wirkt er? Löse die Ungleichung mit dem Rechner (zwei Stellen). \hfill (Abitur 2017 GK) \\ von \leerfeld bis \leerfeld[h]
\end{teile}
\end{aufgabe}
